\documentclass[11pt,letterpaper]{article}
\usepackage[T1]{fontenc}
\usepackage[utf8]{inputenc}
\usepackage{lmodern,microtype,amsmath,amssymb,amsthm,mathtools}
\usepackage[margin=.94in,headheight=14pt]{geometry}
\usepackage{booktabs,tabularx,array,enumitem,xcolor,listings,fancyhdr,xurl,needspace}
\usepackage[colorlinks=true,linkcolor=blue!45!black,urlcolor=blue!45!black,citecolor=blue!45!black]{hyperref}
\definecolor{ink}{HTML}{183B4D}
\definecolor{pale}{HTML}{F2F5F7}
\pagestyle{fancy}\fancyhf{}
\fancyhead[L]{\small Five rules for iPRA}
\fancyhead[R]{\small Canonical recursion and clauses}
\fancyfoot[C]{\thepage}
\newtheorem{theorem}{Theorem}[section]
\newtheorem{lemma}[theorem]{Lemma}
\newtheorem{proposition}[theorem]{Proposition}
\newtheorem{corollary}[theorem]{Corollary}
\theoremstyle{definition}\newtheorem{definition}[theorem]{Definition}
\newcommand{\ipra}{\mathsf{iPRA}}
\newcommand{\Kfive}{\mathcal K_5}
\newcommand{\Ksix}{\mathcal K_6}
\newcommand{\file}[1]{\texttt{\small\nolinkurl{#1}}}
\newcommand{\lit}{\mathrel{\equiv_{\mathrm{syn}}}}
\hypersetup{pdftitle={Five rules for intuitionistic primitive recursive arithmetic},pdfauthor={Research report prepared for Vladimir Reshetnikov},pdfsubject={Canonical recursion, context clauses, and a complete proof checker}}
\lstset{language=Python,basicstyle=\ttfamily\scriptsize,breaklines=true,columns=fullflexible,keepspaces=true,showstringspaces=false,backgroundcolor=\color{pale},frame=single,rulecolor=\color{black!15},xleftmargin=3pt,xrightmargin=3pt}
\setlist{topsep=5pt,itemsep=3pt}
\setlength{\emergencystretch}{2em}
\title{\textcolor{ink}{Five rules for intuitionistic\\primitive recursive arithmetic}\\[4pt]\large Canonical recursion, context clauses, and one opening operation}
\author{Research report prepared for Vladimir Reshetnikov}
\date{11 November 2026}
\begin{document}
\maketitle
\begin{abstract}
We give a smaller presentation of full first-order intuitionistic primitive recursive arithmetic by combining the current canonical successor-seeded recursor with one-term equality transport and a clause presentation of the paired binder. Public formulas use only equality and one two-variable binder; five inference schemes suffice. The new calculus is explicitly equivalent to the repository's six-rule canonical calculus, with both formula composites proved and numerical translations equal to the identity. Its ordinary realization is therefore equiconsistent and bi-interpretable with iPRA on the whole numerical domain, with actual equality and composite graph $x=y$.

A complete, strictly validated standalone Python checker occupies 3,305 normalized bytes, 946 significant tokens, and 1,072 AST nodes. A control implementing the six-rule calculus with the same variable representation occupies 3,820 bytes, 1,107 tokens, and 1,266 nodes. Gzip level 9 sizes are 1,073 and 1,136 bytes. The simplification is a composition of established mathematical ingredients, coupled to an unambiguous tree grammar; it is not a new arithmetic breakthrough or a claim of absolute minimality. The smaller historical compressed scalar checkers remain smaller in raw source size.
\end{abstract}
\setcounter{tocdepth}{1}
\tableofcontents
\newpage
\input{paper/result.tex}
\input{paper/five_rules.tex}
\input{paper/interpretations.tex}
\input{paper/opening.tex}
\input{paper/implementation.tex}
\input{paper/research.tex}
\appendix
\section{The complete primary verifier}
Only the module below and Python's standard JSON/runtime facilities are trusted by its command-line interface. The certificate generators, translators, measurements, and tests are independent supporting programs; the verifier does not import them. Source metrics remove the module docstring and normalize the whole remaining module with \file{ast.unparse}, followed by one line feed. The delivered readable source has the identical acceptance logic.
\lstinputlisting{code/arity_checker.py}
\begingroup\small
\section{Source map and inherited results}
The snapshot is ProveIt commit
\begin{center}\small\texttt{11be5aabc6b07262046e2553419bc6b05271d178}.\end{center}
All repository paths below are relative to
\file{Logic/PrimitiveRecursiveArithmetic/ConciseIPRA/}.
\begin{enumerate}
\item \file{README.md}: current inventory, interpretation standard, and formalization status.
\item \file{Research/canonical_recursors/theory/canonical_R.tex}: canonical substitution coherence, raw/canonical equivalence, six-rule equivalence, and composition with iPRA. The frozen copy is \file{reference/canonical_R.tex}.
\item \file{Research/canonical_recursors/theory/successor_seed_prelogic.tex} and \file{successor_seed_bootstrap.tex}: finite native case repair and the predecessor/initializer recovery. Copies have the same basenames in \file{reference/}.
\item \file{Research/certified_zero/reference/inherited_choice_basis.tex}: the inherited clause interface, logical macros, and formula composites.
\item \file{Research/certified_zero/reference/inherited_ordinary_bridge.tex}, \file{scalar_predecessor.tex}, \file{equality_reflection.tex}, and \file{primitive_recursion.tex}: the ordinary scalar bridge, equality reflection, quantifier-free induction, pairing, and primitive-recursion compiler. Copies are in \file{reference/}.
\item \file{Lean/ConciseIPRA/NativeOneContext.lean} and \file{NativeMinimal.lean}: upstream one-term transport and the combined witness-directed calculus. They concern the upstream scalar kernel, not a formalization of this report.
\end{enumerate}
The archive preserves these textual references with the upstream MIT No Attribution license. \file{evidence/source_provenance.json} records local hashes and upstream Git blob identities. The new theorem is the full five-rule interface and its comparison, including the concrete scope-preserving representation. No priority is claimed for any inherited ingredient.
\begin{thebibliography}{9}
\bibitem{current} ProveIt Contributors, \emph{Concise presentations of intuitionistic primitive recursive arithmetic}, pinned repository snapshot above. \url{https://github.com/VladimirReshetnikov/ProveIt/tree/11be5aabc6b07262046e2553419bc6b05271d178/Logic/PrimitiveRecursiveArithmetic/ConciseIPRA}.
\bibitem{canonical} ProveIt Contributors, \emph{Canonical successor-seeded syntax: a six-rule presentation, with its normalization cost}, \file{Research/canonical_recursors/theory/canonical_R.tex}, same snapshot.
\bibitem{clauses} ProveIt Contributors, \emph{Complementary introductions with one local assumption}, \file{Research/certified_zero/reference/inherited_choice_basis.tex}, same snapshot.
\bibitem{one} ProveIt Contributors, \file{Lean/ConciseIPRA/NativeOneContext.lean}, \file{NativeMinimal.lean}, and the \file{one_context} / \file{one_sided_transport} research packages, same snapshot.
\bibitem{bridge} ProveIt Contributors, \emph{A successor-seeded scalar basis with two computation schemas}; \emph{Pairing, primitive recursion, and internal numerical composites}; source locations in the preceding appendix, same snapshot.
\end{thebibliography}
\endgroup
\end{document}
